\documentclass[dvisvgm,tikz,border=3pt]{standalone}
\usepackage{amsmath,amssymb}
\usepackage{pgfplots}
\usepackage{CJKutf8}
\pgfplotsset{compat=1.18}
\usetikzlibrary{arrows.meta,calc,positioning}

\definecolor{themebg}{HTML}{30362d}
\definecolor{themefg}{HTML}{edf4e8}
\definecolor{themegray}{HTML}{9ea897}
\definecolor{themecurve}{HTML}{7eb6ff}
\definecolor{themealert}{HTML}{ff7b7b}
\definecolor{themeamber}{HTML}{f5b942}

\pgfdeclarelayer{background}
\pgfdeclarelayer{foreground}
\pgfsetlayers{background,main,foreground}

\begin{document}
\begin{CJK*}{UTF8}{gbsn}
\pagecolor{themebg}
\begin{tikzpicture}[color=themefg, text=themefg, >=Stealth]

% ── 几何投影分解图 ──
\begin{scope}[shift={(-1.2,0)}]
  \coordinate (O) at (0,0);
  \coordinate (V1) at (2.6, 0);
  \coordinate (V2) at (2.6, 2.3);
  \coordinate (P) at (2.6, 0);
  \coordinate (U2) at (0, 2.3);

  % 坐标原点
  \node[below left, font=\small] at (O) {$O$};

  % 第一方向 u1 = v1 (晶蓝)
  \draw[->, themecurve, line width=2.4pt] (O) -- (V1) node[below=2pt, font=\small] {$\boldsymbol{u}_1 = \boldsymbol{v}_1$};

  % 原始第二向量 v2 (中性灰)
  \draw[->, themegray, line width=2.0pt] (O) -- (V2) node[above left, font=\small] {原向量 $\boldsymbol{v}_2$};

  % 投影向量 proj_{u1}(v2) (暖金粗实线)
  \draw[->, themeamber, line width=2.0pt] (O) -- ($(O)!0.95!(P)$) node[midway, below=3pt, font=\footnotesize, text=themeamber] {$\operatorname{proj}_{\boldsymbol{u}_1}\boldsymbol{v}_2$};

  % 正交补充向量 u2 = v2 - proj(v2) (珊瑚红粗实线，标签置于左侧三角形内)
  \draw[->, themealert, line width=2.4pt] (P) -- (V2) node[midway, fill=themebg, inner sep=1pt, left=2pt, font=\scriptsize, text=themealert] {$\boldsymbol{u}_2 = \boldsymbol{v}_2 - \operatorname{proj}_{\boldsymbol{u}_1}\boldsymbol{v}_2$};

  % 原点平移后的正交基 u2
  \draw[->, themealert, line width=2.0pt] (O) -- (U2) node[left=2pt, font=\small, text=themealert] {正交基 $\boldsymbol{u}_2 \perp \boldsymbol{u}_1$};

  % 直角标记
  \draw[themegray, line width=0.8pt] ($(P)+(-0.25, 0)$) -- ($(P)+(-0.25, 0.25)$) -- ($(P)+(0, 0.25)$);
  \draw[themegray, line width=0.8pt] ($(O)+(0.25, 0)$) -- ($(O)+(0.25, 0.25)$) -- ($(O)+(0, 0.25)$);
\end{scope}

% ── 右侧：代数契约与通式说明 ──
\begin{scope}[shift={(6.2, 1.3)}]
  \node[draw=themegray!60, fill=themebg, rounded corners=6pt, line width=0.8pt, inner sep=8pt, text width=6.6cm] {
    \centering
    {\bfseries\color{themecurve} Gram--Schmidt 核心代数契约} \\[6pt]
    \raggedright\small
    $\bullet$ \textbf{第一步}：$\boldsymbol{u}_1 = \boldsymbol{v}_1$（固定第一个方向）\\
    $\bullet$ \textbf{第二步}：$\boldsymbol{u}_2 = \boldsymbol{v}_2 - \frac{\langle\boldsymbol{v}_2,\boldsymbol{u}_1\rangle}{\langle\boldsymbol{u}_1,\boldsymbol{u}_1\rangle}\boldsymbol{u}_1$（剔除已有方向分量）\\
    $\bullet$ \textbf{第 $j$ 步通式}：$\boldsymbol{u}_j = \boldsymbol{v}_j - \sum_{i=1}^{j-1}\frac{\langle\boldsymbol{v}_j,\boldsymbol{u}_i\rangle}{\langle\boldsymbol{u}_i,\boldsymbol{u}_i\rangle}\boldsymbol{u}_i$\\
    $\bullet$ \textbf{规范化（最后统一）}：$\boldsymbol{q}_i = \frac{\boldsymbol{u}_i}{\|\boldsymbol{u}_i\|}$\\[4pt]
    {\color{themealert}\bfseries 核心不变量}：$\operatorname{span}\{\boldsymbol{u}_1,\ldots,\boldsymbol{u}_j\} = \operatorname{span}\{\boldsymbol{v}_1,\ldots,\boldsymbol{v}_j\}$
  };
\end{scope}

% ── 底部总结横条 ──
\node[font=\bfseries\small, color=themecurve, anchor=north] at (4.8, -1.6) {
  核心心智模型：正交化是在保持Span 前提下消除方向耦合；“减投影”就是剪掉已有的冗余相关分量
};

\end{tikzpicture}
\end{CJK*}
\end{document}
